\begin{qfresult}[The Campbell--Shiller log-linear present-value identity: the exact return identity, its first-order approximation and the constant $\rho$, the forward solution, and why a valuation ratio must forecast returns, dividend growth or itself]\label{res:campbell-shiller-identity}
\qffield{Statement} \emph{The exact identity.} With $R_{t+1}=(P_{t+1}+D_{t+1})/P_t$ and lowercase letters for logs, so that $pd_t=p_t-d_t$,
\[ r_{t+1}=\Delta d_{t+1}-pd_t+\log\bigl(1+e^{pd_{t+1}}\bigr). \]
It is the definition of a return rearranged and applies to each data point \citep[RFS 21(4), 2008, p.~1539 and n.~4]{wide-cochrane-dog}; no model of prices is involved.

\emph{The approximation.} Expanding $\log(1+e^{x})$ to first order around the mean $\overline{pd}$ gives
\[ r_{t+1}\approx k+\rho\,pd_{t+1}-pd_t+\Delta d_{t+1},\qquad \rho=\frac{e^{\overline{pd}}}{1+e^{\overline{pd}}}, \]
with $k=\log(1+e^{\overline{pd}})-\rho\,\overline{pd}$. Cochrane writes it for demeaned variables, with $\rho=PD/(1+PD)$ at the price--dividend ratio about which one linearises. In annual CRSP data for 1927--2004 he takes $\rho$ from the mean log dividend yield, $\rho=0.9638$ (eq.~(4), pp.~1539--1540). The constant $k$ is our own bookkeeping (Interpretation).

\emph{Forward solution.} Solving the approximate identity forward $J$ periods gives
\[ pd_t=\sum_{j=1}^{J}\rho^{j-1}\bigl(k+\Delta d_{t+j}-r_{t+j}\bigr)+\rho^{J}pd_{t+J}, \]
apart from the approximation errors. It holds after the event, and therefore also in expectation. If $\rho^{J}pd_{t+J}$ vanishes as $J$ grows (no bubble), a high price--dividend ratio must be followed by high dividend growth or low returns. Cochrane's eq.~(9) is the expectational form, $d_t-p_t=E_t\sum_j\rho^{j-1}r_{t+j}-E_t\sum_j\rho^{j-1}\Delta d_{t+j}$ (p.~1544).

\emph{Restrictions on forecasts.} Project each term of the identity on $d_t-p_t$. With $b_r$, $b_d$ and $\phi$ the slopes of next period's log return, dividend growth and dividend yield on $d_t-p_t$,
\[ b_r=1-\rho\phi+b_d \]
(eq.~(5), p.~1540). While $\phi<1/\rho\approx1.04$, returns and dividend growth cannot both be unforecastable. A coherent null of no return predictability therefore needs $b_d<0$, a forecast of dividend growth that can itself be tested. Dividing by $1-\rho\phi$ gives long-run coefficients that sum to one, $b_r^{lr}-b_d^{lr}=1$, the shares of the dividend yield's variance due to expected returns and expected dividend growth (eqs.~(8)--(11), p.~1544).

\qffield{Conditions} \emph{Two versions.} The working paper linearises the present-value model itself around a constant discount rate $r$ and growth $g$, with $\rho=\exp(-(r-g))$ \citep[NBER WP 2100, eqs.~(2)--(5), pp.~4--5]{campbell-shiller-1986-nber}. Its return $R_t=(P_{t+1}-P_t+D_t)/P_t$ dates the price at the start of period $t$, and its log dividend--price ratio $\delta_t=d_{t-1}-p_t$ divides the dividend of the period just ended by that price. With end-of-period prices and trailing dividends, as in the data below, this is the timing used here (Interpretation). The same model follows from linearising the one-period return around $\Delta d_t=g$ and $\delta_t=\log(e^{r-g}-1)$. For $r$ the paper takes the log of one plus the sample average real return, $\log(1+\bar R)$, not the average log return, and for $g$ the average change in log dividends (p.~5). In a steady state the two constants agree, since $e^{\bar\delta}=e^{r-g}-1$ makes $\rho=1/(1+e^{\bar\delta})$ (Interpretation). In a finite sample they differ. For nominal annual S\&P~500 data 1872--2025, $\rho$ is 0.9634 from the mean log price--dividend ratio and 0.9364 from $\log(1+\bar R)-\bar g$. Two effects lower the second: $\log(1+\bar R)$ exceeds the mean log return, and the ratio rose over the sample, which adds its average change to the mean log return (Interpretation).

\emph{Accuracy.} For annual PPI-deflated returns 1871--1986, at the working paper's point, the error variance of the one-period approximation is 0.0008 of the variance of the log return, with a correlation of 0.9996 \citep[Table~1a, p.~21]{campbell-shiller-1986-nber}. The approximation degrades when the ratio moves far from the point of linearisation. Cochrane notes that his eq.~(5) is not exact: each coefficient differs from the one implied by the other two by about 0.005 \citep[p.~1540]{wide-cochrane-dog}.

\emph{Units.} Inflation cancels from $r_{t+1}-\Delta d_{t+1}$, so real and nominal versions share $pd_t$ and, at the same expansion point, the same approximation error; only the variance of $r$ differs. The working paper's point does change with the deflator, through $\log(1+\bar R)$ and $\bar g$. Subtracting the risk-free rate from both sides turns ``dividend growth'' into dividend growth less the interest rate (p.~1540). Cash dividends omit repurchases, which matters for forecasting with the dividend yield (p.~1538).

\emph{Finite horizons.} The terminal term does not vanish quickly: with $\rho=0.963$ its weight $\rho^J$ is 0.83 after 5 years and 0.47 after 20.

\qffield{Practice} \emph{S\&P~500, annual} (Goyal's file of 2026: index level and trailing twelve-month dividends). The return is built from the same two series in every year, $R_t=(\text{Index}_t+\text{D12}_t)/\text{Index}_{t-1}-1$, and the exact identity then holds to rounding error ($2\times10^{-15}$) for 1872--2025. CRSP's S\&P~500 return, used from 1926 in the forecasting examples of the next two entries, is a different series. Its log return differs from this one by up to 6.8 log points (1931), with a standard deviation of 1.0 point over 1926--2025, so it is kept out of this check. The approximation error is then exactly the Taylor remainder $\log(1+e^{pd})-\log(1+e^{\overline{pd}})-\rho\,(pd-\overline{pd})$. It is never negative, because $\log(1+e^{x})$ is convex, and never above $(pd-\overline{pd})^2/8$, because the second derivative is at most $1/4$. Nominal returns:
\begin{center}\small
\begin{tabular}{@{}llrrr@{}}
\toprule
Sample & Expansion point & $\rho$ & Error variance / variance of $r$ & Correlation \\
\midrule
1872--1986 & mean $pd$ & 0.9547 & 0.00021 & 0.99991 \\
 & $\log(1+\bar R)-\bar g$ & 0.9348 & 0.00089 & 0.99960 \\
1872--2025 & mean $pd$ & 0.9634 & 0.00077 & 0.99965 \\
 & $\log(1+\bar R)-\bar g$ & 0.9364 & 0.00728 & 0.99638 \\
\bottomrule
\end{tabular}
\end{center}
This is an illustration on nominal data, not a reproduction of Campbell and Shiller's Table~1a: their 0.0008 and 0.9996 are for PPI-deflated returns, and the deflator is not in the file used here. The working paper's point, $\overline{pd}=2.69$ for 1872--2025, lies far below the ratio's recent values, so its error grows once the later decades are included. If CRSP returns replaced the constructed ones from 1926, the mean-$pd$ error ratios would be 0.0023 and 0.0025, mostly the mismatch between the two return series rather than the approximation.

\emph{Decomposition, 1927--2004} (outcome years, predictor lagged one year). Slopes on $d_{t-1}-p_{t-1}$: $b_r=0.066$, $b_d=-0.067$, $\phi=0.896$, $\rho=0.964$. Then $b_r-b_d+\rho\phi=0.997$, and the gap from one equals the slope of the approximation error, $-0.003$. The long-run shares are 0.49 for returns and 0.49 for dividend growth; for 1872--2025 they are 0.21 and 0.76. With CRSP returns in place of the constructed ones, the 1927--2004 return slope is 0.064.

For the same years, Cochrane's CRSP value-weighted real data give $b_r\approx0.10$, $b_d\approx0$ and $\phi\approx0.94$ \citep[p.~1540]{wide-cochrane-dog}, so almost all the variance falls on returns. The split depends on the index, the dividend series and the sample. These are point estimates with large standard errors, about 0.05 each in Cochrane's sample.

\qffield{Pitfalls} Treating the identity as evidence of predictability: it only says that something must be forecastable, not what. Mixing timing conventions ($d_t$ against $d_{t-1}$ in the ratio) between the identity and the regressions. Checking the identity with a return built from other prices or dividends than the ratio, which adds their mismatch to the ``approximation error''. Computing the working paper's point from the mean log return instead of $\log(1+\bar R)$. Using a $\rho$ estimated in one sample to decompose another. Ignoring the terminal term at horizons of a decade or two. Computing returns from a price index without dividends.

\qffield{See also} \hyperref[meth:gordon-growth-model]{Gordon growth model} (the constant-growth case, in which $pd$ is constant); \hyperref[meth:autoregressive-forecasting]{autoregressive forecasting} (the persistence $\phi$); \hyperref[meth:nested-forecast-tests]{nested forecast tests} (out-of-sample evidence on $b_r$); \hyperref[meth:restricted-forecast-value]{restricted forecasts}.

\end{qfresult}

\begin{qfmethod}[Comparing a predictive regression with the historical mean out of sample: recursive forecasts, out-of-sample $R^2$ and $\Delta$RMSE, MSE-F and ENC, the upward bias of the larger model's MSPE, and the Clark--West adjusted test]\label{meth:nested-forecast-tests}
\qffield{Statement} \emph{Design.} Model~1, the historical mean, is nested in model~2, a regression of the return on a lagged predictor. With $R$ initial observations and $P$ forecasts, each forecast for period $t+1$ uses estimates from data through $t$ only. In the \emph{recursive} scheme the estimation sample grows; in the \emph{rolling} scheme it keeps length $R$ \citep[NBER Technical WP 326, sec.~3, eq.~(3.6), p.~6, and n.~3]{clark-west-2006-nber}. Welch and Goyal's forecasts likewise use ``only the data available up to the time at which the forecast is made'' \citep[p.~1461]{s7-welch}.

\emph{Statistics.} With $e_N$ and $e_A$ the errors of the mean and the regression over the $P$ forecasts:
\begin{itemize}
\item the out-of-sample $R^2_{OS}=1-\sum(r_t-\hat r_t)^2/\sum(r_t-\bar r_t)^2$, comparable in units with the in-sample $R^2$ and of the same sign as $\Delta\mathrm{RMSE}=\sqrt{\mathrm{MSE}_N}-\sqrt{\mathrm{MSE}_A}$ \citep[NBER WP 11468, eq.~(1), p.~6]{campbell-thompson-2005-nber};
\item McCracken's $\text{MSE-F}=(T-h+1)(\mathrm{MSE}_N-\mathrm{MSE}_A)/\mathrm{MSE}_A$, a test of equal MSE \citep[eq.~(6), p.~1461]{s7-welch};
\item the encompassing statistic $\text{ENC}=\frac{T-h+1}{T}\sum_t(e_{Nt}^2-e_{Nt}e_{At})/\mathrm{MSE}_A$ (eq.~(7), p.~1462).
\end{itemize}
Here $T$ is the number of forecasts and $h$ the overlap ($h=1$ for none). For nested models both MSE-F and ENC have nonstandard distributions: under the null the asymptotic difference in squared errors is exactly zero with zero variance. Welch and Goyal therefore bootstrap their critical values (p.~1462).

\emph{Why the larger model looks worse.} Under the null the extra slope is zero in population, but the regression estimates it and adds noise to its forecasts. With $\hat y_1$ and $\hat y_2$ the two forecasts,
\[ \hat\sigma_1^2-\hat\sigma_2^2=-2P^{-1}\textstyle\sum\hat e_{1}(\hat y_{1}-\hat y_{2})-P^{-1}\sum(\hat y_{1}-\hat y_{2})^2 . \]
The first term should be near zero under the null, and the second is negative. So the parsimonious model is expected to have the smaller sample MSPE even when the null is true \citep[eq.~(3.8), p.~7]{clark-west-2006-nber}.

\emph{The Clark--West test.} Add back the noise term and test whether the adjusted difference is positive. Define
\[ \hat f_{t+\tau}=(y_{t+\tau}-\hat y_{1t,t+\tau})^2-\bigl[(y_{t+\tau}-\hat y_{2t,t+\tau})^2-(\hat y_{1t,t+\tau}-\hat y_{2t,t+\tau})^2\bigr], \]
regress it on a constant, and reject equal MSPE in favour of the larger model if the $t$ statistic exceeds 1.282 (one-sided 10 percent) or 1.645 (5 percent). For one-step forecasts the ordinary standard error serves; autocorrelated multistep errors need an autocorrelation-consistent one (eq.~(2.1), p.~4).

\emph{Approximate size.} The statistic is not asymptotically normal. Clark and West rely on Clark and McCracken's asymptotics, in which $R\to\infty$ and $P\to\infty$ with $P/R$ tending to a finite nonzero constant. For one-step forecasts with conditionally homoskedastic errors, the limiting distribution then depends only on the number of extra regressors, the limit of $P/R$ and the scheme. Clark and McCracken tabulated its quantiles numerically for up to 20 extra regressors and $P/R$ from 0.1 to 20. In those tables the 0.90 quantile lies at or below 1.282 and the 0.95 quantile at or above it, apart from a couple of exceptions. All 400 recursive entries obey both bounds, but in 4 of the 400 rolling entries the 0.90 quantile is 1.29. Within that regime a nominal 0.10 test therefore has an actual size somewhere between 0.05 and 0.10 (eq.~(4.2), pp.~8--10). The authors call this safe in practice but stress that it is not proved, that they do not assert it is literally true, and that further exceptions may appear. With 1.645 the corresponding range is 0.01 to 0.05, with one recursive and 14 rolling exceptions (eq.~(4.3), p.~10). With one extra regressor the same quantiles also serve for multistep and conditionally heteroskedastic forecasts; otherwise Clark and McCracken simulate critical values (pp.~10--12).

The mean of $\hat f$ equals $2P^{-1}\sum\hat e_1(\hat e_1-\hat e_2)$, an encompassing form (eq.~(4.1), p.~8). With $h=1$, $\text{ENC}=P\,\bar f/(2\,\mathrm{MSE}_A)$, so ENC and Clark--West share a numerator (Interpretation).

\qffield{Conditions} \emph{Scope of the theory.} Nested linear models estimated by least squares, with multistep forecasts made directly rather than iterated \citep[sec.~4, p.~8]{clark-west-2006-nber}.

\emph{Simulation evidence, not a guarantee.} Clark and West ran 48 one-step designs: two DGPs, homoskedastic and conditionally heteroskedastic errors, rolling and recursive schemes, and several $R$ and $P$ (Tables~1--2). In every design the nominal 0.10 test had a size between 0.05 and 0.10, with a median of 0.080. At 1.645 the sizes ranged from 0.027 to 0.059, outside 0.01--0.05 in three designs. Rolling regressions gave sizes closer to the nominal level than recursive ones, and the unadjusted comparison with normal critical values had a median size of 0.008 (pp.~2, 15--17). This is evidence for their designs, not a bound for every DGP, sample size or forecast horizon.

\emph{What is tested.} The test asks whether the predictor's population slope is non-zero, not whether the estimated regression beat the mean in this sample. A negative $\Delta$RMSE can therefore coexist with a significant test: these tests ``are ultimately tests of whether $\beta$ is equal to zero'' \citep[p.~1462, n.~3]{s7-welch}. Cochrane adds that out-of-sample $R^2$ ``is not a \emph{test}'' with more power than full-sample tests \citep[p.~1538]{wide-cochrane-dog}.

\emph{Adjusted $R^2$.} Welch and Goyal's tables report an adjusted out-of-sample $\bar R^2$. The formula printed with eq.~(6), $\bar R^2=R^2-(1-R^2)(T-k)/(T-1)$, would put an $R^2$ near zero at about $-100$ percent, which the tables do not show; it is presumably misprinted (Interpretation). $\Delta$RMSE needs no adjustment and is the safer basis for comparison.

\emph{Data.} The examples below use Goyal's data file as served in 2026, which runs to 2025; the 2008 tables used data to 2005, and revisions may move the numbers.

\qffield{Practice} \emph{Size and power by simulation.} Clark and West's DGP~1 is calibrated to monthly S\&P excess returns (percent) and the dividend--price ratio: $y_t=0.5+\gamma z_{t-1}+e_t$, $z_t=0.15+0.95z_{t-1}+v_t$, $\mathrm{var}(e)=18$, $\mathrm{var}(v)=0.025$, $\mathrm{corr}(e,v)=-0.75$ \citep[eq.~(6.1), p.~13]{clark-west-2006-nber}. With the recursive scheme, $R=120$ and 5,000 replications, the rejection rates at 1.282 are below, with Clark and West's Table~1, panel~B1, in brackets (unnumbered table page). Under the null, the mean of $\mathrm{MSE}_N-\mathrm{MSE}_A$ is $-0.160$ ($P=120$) and $-0.118$ ($P=240$): the historical mean wins on average.
\begin{center}\small
\begin{tabular}{@{}lrr@{}}
\toprule
Rejection rate, nominal 0.10 & $P=120$ & $P=240$ \\
\midrule
Clark--West, $\gamma=0$ & 0.069 (0.085) & 0.059 (0.070) \\
Unadjusted MSPE difference, $\gamma=0$ & 0.023 (0.028) & 0.010 (0.012) \\
Clark--West, $\gamma=0.35$ & & 0.156 \\
Unadjusted MSPE difference, $\gamma=0.35$ & & 0.033 \\
\bottomrule
\end{tabular}
\end{center}
These Clark--West sizes lie in the predicted band but below the published ones. The gaps are 0.016 and 0.011, or 3.0 and 2.3 standard errors once the simulation error of both studies (5,000 replications each) is counted. A rerun in Clark and West's own layout used one path per replication shared by every $P$, with no observed lag for the first observation, and 20,000 replications (p.~15). It gives 0.072, 0.066 and 0.055 for $P=120$, 240 and 720, against 0.085, 0.070 and 0.054. For the rolling scheme with $P=720$, it reproduces the means they quote, $-0.24$, 0.01 and 0.25, as $-0.247$, 0.005 and 0.252 (p.~16). The remaining gap is at $P=120$: about 0.013, or three standard errors, in both schemes. Its cause is unresolved. The $\gamma=0.35$ rows are raw rejection rates, not size-adjusted power.

\emph{Welch and Goyal's dividend--price row.} The annual log equity premium is regressed on the lagged log dividend--price ratio, 1872--2005, with the first forecast for 1892 after 20 observations. The result is $R^2_{OS}=-1.32$ percent and $\Delta\mathrm{RMSE}=-0.13$ points against their $-0.11$ \citep[Table~1, p.~1470]{s7-welch}, MSE-F $-1.49$, ENC 0.32 and a Clark--West $t$ of 0.26. The conventionally adjusted value is $-2.23$ against their $-2.06$. From 1965 the figures are $\Delta\mathrm{RMSE}=-0.13$ (theirs $-0.09$), $R^2_{OS}=-1.67$ percent and a Clark--West $t$ of 0.46. The in-sample adjusted $R^2$ is 0.32 percent against their 0.49.

\emph{Start of the evaluation sample.} Same regression, recursive from 1872:
\begin{center}\small
\begin{tabular}{@{}lrrrr@{}}
\toprule
& \multicolumn{2}{c}{Ending 2005} & \multicolumn{2}{c}{Ending 2025} \\
First forecast & $R^2_{OS}$ (\%) & Clark--West $t$ & $R^2_{OS}$ (\%) & Clark--West $t$ \\
\midrule
1892 & $-1.32$ & 0.26 & $-2.26$ & $-0.11$ \\
1927 & $-0.87$ & 0.48 & $-2.13$ & 0.06 \\
1947 & $-0.93$ & 0.63 & $-3.10$ & 0.17 \\
1965 & $-1.67$ & 0.46 & $-4.23$ & $-0.01$ \\
1985 & $-17.40$ & $-0.65$ & $-12.18$ & $-0.98$ \\
\bottomrule
\end{tabular}
\end{center}
No start gives a positive $R^2_{OS}$ or a Clark--West $t$ near 1.282, and adding 2006--2025 makes every row worse. The 1985 rows rest on 21 and 41 forecasts.

\emph{An adjusted test against a raw loss.} In Clark and West's monthly example (S\&P excess return on the dividend--price ratio, recursive, forecasts 1964--2004, $P=492$), the regression's MSPE is 19.14 against 18.91 for the mean. The adjustment is 0.68, so the adjusted MSPE is 18.46, and the Clark--West $t$ is 1.17. That is below 1.282 but above Clark and McCracken's tabulated 0.90 quantile, which Clark and West read as ``some statistical evidence'' against the null (pp.~19--20; the results table is printed as Table~4 and called Table~7 in the text).

\qffield{Pitfalls} Applying the Diebold--Mariano or another normal test to the raw MSPE difference of nested models: it is undersized and almost never rejects. Reading a negative $R^2_{OS}$ as proof that the slope is zero, or a positive one as proof of a usable strategy. Choosing the start of the evaluation sample after seeing the results. Two-sided tests: the alternative is that the larger model forecasts better. Overlapping multi-year forecasts with ordinary standard errors (\hyperref[meth:correlation-independence-tests]{overlapping observations}). Fixing a coefficient's expected sign from the full sample, which leaks the future into a ``real-time'' forecast (\hyperref[meth:restricted-forecast-value]{restricted forecasts}).

\qffield{See also} \hyperref[meth:forecast-accuracy-evaluation]{chronological forecast evaluation} (the split and the benchmarks); \hyperref[meth:hypothesis-testing]{hypothesis tests} (size and power); \hyperref[res:campbell-shiller-identity]{present-value identity} (why the dividend yield should forecast); \hyperref[meth:backtest-evaluation-diagnostics]{backtest diagnostics} (Welch and Goyal's other predictors); \hyperref[meth:forecast-combination]{forecast combination}; \hyperref[meth:conditional-predictive-ability]{conditional predictive ability} (comparing forecasting methods with fixed windows).

\end{qfmethod}

\begin{qfmethod}[Restricted equity-premium forecasts and their economic value: sign restrictions on the slope and on the forecast, the mean--variance weight with no short sales and limited leverage, and the utility gain over the historical mean]\label{meth:restricted-forecast-value}
\qffield{Statement} \emph{Restrictions.} A regression estimated on a short early sample can produce perverse forecasts. Campbell and Thompson's example is the early 1930s: the earnings--price ratio was high, but its estimated coefficient was negative, so the regression forecast a negative premium \citep[NBER WP 11468, p.~6]{campbell-thompson-2005-nber}. They impose two restrictions, separately and together (pp.~2--3, 6):
\begin{itemize}
\item set the slope to zero when it has the ``wrong'' sign, which leaves the historical mean as the forecast;
\item set a negative forecast of the equity premium to zero, and floor the historical-mean benchmark in the same way.
\end{itemize}

\emph{Economic value.} An investor with mean--variance preferences and relative risk aversion $\gamma$ who observes the predictor holds the share $\alpha_t=(1/\gamma)(\mu+x_t)/\sigma_\varepsilon^2$ in stocks (eq.~(5), p.~9). In Campbell and Thompson's portfolio exercise, $\gamma=3$ and the weight is confined to 0--150 percent: no short sales and at most 50 percent leverage. The variance is estimated either from all data up to the investment date or from a rolling five-year window (p.~10). Each forecast is scored by the average realised utility
\[ U=\widehat E(R_p)-\tfrac{\gamma}{2}\widehat{\mathrm{Var}}(R_p),\qquad R_{p,t}=R_{f,t}+w_t\,(R_t-R_{f,t}) \]
\citep[Table~2 notes, unnumbered table page]{campbell-thompson-2005-nber}. The predictor's value is the change in $U$ from using it instead of the historical mean. Utility differences are in units of expected return, so they read as the monthly fee an investor would pay for the predictor (p.~10). How a small $R^2$ translates into a proportional return gain, about $R^2/S^2$, is set out under \hyperref[meth:backtest-evaluation-diagnostics]{backtest diagnostics}.

\qffield{Notation} $\hat\mu_t$ is the forecast of next month's excess return and $\hat\sigma_t^2$ the variance of excess returns observed before month $t$; the bounded weight is
\[ w_t=\min\bigl\{\max\{\hat\mu_t/(\gamma\hat\sigma_t^2),\,0\},\,1.5\bigr\}. \]
Unconstrained, this is the optimal risky share $y^\ast$ of the \hyperref[meth:capital-allocation-line]{capital allocation line}, with $A=\gamma$ for returns in decimals (\hyperref[def:risk-aversion-utility]{risk aversion and utility}).

\qffield{Conditions} \emph{Real time.} Campbell and Thompson define the ``wrong'' sign as one ``different from the theoretically expected sign estimated over the full sample'' (p.~6). Taking the expected sign from the full-sample estimate uses information not available at the forecast date. Fixing it in advance from theory, positive for valuation ratios, avoids this (Interpretation).

\emph{Choices that move the answer.} The value of $\gamma$, the weight bounds, the variance estimator, monthly rebalancing and zero trading costs. The utility is computed from realised means and variances over one evaluation sample and carries no standard error.

\emph{Version.} Locators refer to the working paper of June 2005; the \emph{Review of Financial Studies} version (2008) was not used here and may differ.

\qffield{Practice} \emph{Hand example}, $\gamma=3$, $R_f=0.3$ percent a month:
\begin{center}\small
\begin{tabular}{@{}rrrrr@{}}
\toprule
$\hat\mu_t$ & $\hat\sigma_t^2$ & $\hat\mu_t/(\gamma\hat\sigma_t^2)$ & $w_t$ & $R_{p,t}$ when $R_t=2.0,\,-5.0,\,1.0\%$ \\
\midrule
0.6\% & 0.0020 & 1.00 & 1.00 & 2.00\% \\
$-0.2$\% & 0.0020 & $-0.33$ & 0 & 0.30\% \\
3.0\% & 0.0018 & 5.56 & 1.50 & 1.35\% \\
\bottomrule
\end{tabular}
\end{center}
The bounds bind twice. With the sample variance of the three portfolio returns, $U=0.012167-1.5\times0.0000736=0.012056$, or 1.206 percent a month.

\emph{Campbell and Thompson's dividend--price row.} Monthly S\&P~500, regression data from February 1871, forecasts January 1927 to December 2003, $\gamma=3$, weights in [0, 1.5], variance from all earlier excess returns, expected sign positive:
\begin{center}\small
\begin{tabular}{@{}lrr@{}}
\toprule
& Recomputed & Table~2 \\
\midrule
$R^2_{OS}$, restricted (\%) & 0.003 & 0.001 \\
$R^2_{OS}$, unrestricted (\%) & $-0.326$ & \\
Utility, historical mean (\% a month) & 0.492 & 0.501 \\
Change, unrestricted (\% a month) & $-0.018$ & $-0.016$ \\
Change, restricted (\% a month) & $-0.024$ & $-0.024$ \\
\bottomrule
\end{tabular}
\end{center}
Published figures are from Table~2 (constant-variance columns) and its panel~D source \citep[Table~2, unnumbered table page]{campbell-thompson-2005-nber}. The restrictions lift $R^2_{OS}$ from $-0.33$ to about zero but do not improve the utility change, which moves from $-0.018$ to $-0.024$. The slope was negative in 8 percent of months, and the no-short-sale bound already removes negative forecasts from the portfolio, as Campbell and Thompson observe (p.~10). Extended to December 2025, both changes fall to about $-0.05$ percent a month.

\emph{When the predictor is real.} Clark and West's DGP~1 (see \hyperref[meth:nested-forecast-tests]{nested forecast tests}) is rescaled to monthly decimals with a mean excess return of 0.5 percent and true slope 0.35, a population $R^2$ of 0.17 percent, $S^2=0.014$. With $R=240$ and $P=480$, over 1,000 samples, the mean utility gain is 0.017 percent a month unrestricted and 0.021 percent restricted. With known parameters and no bounds it would be about 0.03 percent, $(1/(2\gamma))(R^2/(1-R^2))(1+S^2)$, half of Campbell and Thompson's eq.~(8) (Interpretation). Only 51 and 54 percent of 40-year samples show a gain.

With DGP~1's original mean of 1.55 percent a month, the optimal weight is near 2.9. The 1.5 bound then binds in 98 to 100 percent of months, and the predictor can hardly be used.

\qffield{Pitfalls} Setting the expected sign from the full sample. Restricting the model's forecast but not the benchmark's. Reading a better $R^2_{OS}$ as a better portfolio when the weight bounds already do the restricting. Mixing percent and decimal returns in $\hat\mu/(\gamma\hat\sigma^2)$, which changes the weight by a factor of 100. Treating a utility gain from one sample as an expected fee without its sampling error.

\qffield{See also} \hyperref[meth:nested-forecast-tests]{nested forecast tests} (the statistical comparison); \hyperref[res:campbell-shiller-identity]{present-value identity} (the expected sign); \hyperref[meth:capital-allocation-line]{capital allocation line}; \hyperref[meth:backtest-evaluation-diagnostics]{backtest diagnostics}.

\end{qfmethod}

\begin{qfmethod}[Combining equity-premium forecasts: equal-weight and performance-weighted averages of single-predictor regressions, why correlated errors limit diversification, and the gain as the spread of the forecasts]\label{meth:forecast-combination}
\qffield{Statement} \emph{Design.} Each of $N$ predictors gets its own regression $r_{t+1}=\alpha+\beta x_{i,t}+\varepsilon_{t+1}$, estimated recursively as in \hyperref[meth:nested-forecast-tests]{nested forecast tests}. The combination forecast is a weighted average of the $N$ individual forecasts,
\[ \hat r_{c,t+1}=\sum_{i=1}^{N}\omega_{i,t}\,\hat r_{i,t+1}, \]
with weights fixed \emph{ex ante} at time $t$ \citep[working paper, February 2008, eqs.~(1)--(2), pp.~3--4]{rapach-strauss-zhou-2008-wp}. The simple schemes are the mean ($\omega_{i,t}=1/N$), the median, and the trimmed mean, which drops the highest and lowest forecasts and weights the others $1/(N-2)$ (pp.~4--5). They need no holdout period.

\emph{Performance weights.} The discounted MSPE scheme weights each forecast by its inverse discounted squared error over a holdout period, $\omega_{i,t}=m_{i,t}^{-1}/\sum_j m_{j,t}^{-1}$ with $m_{i,t}=\sum_{s=R}^{t-1}\theta^{t-1-s}(r_{s+1}-\hat r_{i,s+1})^2$. With $\theta=1$, this is Bates and Granger's optimal combination when forecast errors are unbiased and uncorrelated, so their second-moment matrix is diagonal (Interpretation; the working paper calls the forecasts uncorrelated); $\theta<1$ favours recent accuracy (eqs.~(3)--(4), p.~5). Cluster and approximate Bayesian model-averaging weights complete the list (pp.~5--6).

\emph{Why averaging helps, and where it stops} (Interpretation). With equal weights the combination's error is the average of the individual errors, so
\[ \mathrm{MSE}_c=\frac1{N^2}\sum_{i,j}E[e_ie_j]=\frac1N\sum_i\mathrm{MSE}_i-\frac1N\sum_iE\bigl[(\hat r_i-\hat r_c)^2\bigr]. \]
The second form holds exactly in any sample. The combination is never worse than its average member, and it gains exactly the mean squared spread of the forecasts about their average. It can still lose to the best member. If the errors are unbiased, have equal variance $\sigma^2$ and common correlation $\rho$, then $\mathrm{MSE}_c=\sigma^2[\rho+(1-\rho)/N]$. With individual error biases $b_i$ and the same equal-variance, common-correlation structure for the centered errors, add $\bar b^2$, where $\bar b=N^{-1}\sum_i b_i$. For fixed $0\leq\rho\leq1$, the unbiased-error expression tends to $\rho\sigma^2$ as $N$ grows.

\emph{What is diversified.} Return forecast errors share the unpredictable part of the return, so their correlations are close to one. Write each forecast as the historical mean plus a deviation. The deviations, the predictors' signals, are much less correlated, and averaging them shrinks the combination towards the historical mean when the signals disagree (Interpretation).

\emph{Estimated weights.} Among weights that sum to one, $\Sigma^{-1}\mathbf 1/(\mathbf 1'\Sigma^{-1}\mathbf 1)$ minimises the MSE when $\Sigma$, the matrix of second moments of the errors, is known. With a diagonal $\Sigma$ it reduces to inverse-MSE weights, the $\theta=1$ case above. In practice $\Sigma$ must be estimated from a short holdout, so the weights carry estimation error that equal weights avoid (Interpretation).

\qffield{Conditions} \emph{Version.} Locators refer to the working paper of 6 February 2008, ``Out-of-Sample Equity Premium Prediction: Consistently Beating the Historical Average''. The \emph{Review of Financial Studies} article (2010) has a different title and may differ in design and results; its tables are not quoted here.

\emph{The working paper's evidence.} Monthly data, 14 of Welch and Goyal's predictors available from 1927, evaluation 1947:01--2005:12. Only four unrestricted single-predictor forecasts have a positive $R^2_{OS}$, none above 0.33 percent; the mean combination has 0.95 percent, with a Clark--West $p$-value of 0.09 percent \citep[pp.~12--13, 16 and Table~1, panels A--B, unnumbered table page]{rapach-strauss-zhou-2008-wp}. The authors add a data-mining bootstrap over all predictors and combining methods, because several were tried (sec.~2.5, pp.~11--12). They also warn that the Clark--West test has not been analysed for combination forecasts, which are not nested linear models (n.~11, p.~9).

\emph{Real time.} The working paper takes the expected signs for its restricted forecasts from full-sample estimates (n.~9, p.~8), the leak discussed under \hyperref[meth:restricted-forecast-value]{restricted forecasts}. The predictors, the start of the evaluation and the combining method must likewise be fixed before the evaluation sample is seen.

\qffield{Practice} \emph{Four annual predictors} (Goyal's data file of 2026, a vintage that runs to 2025). The outcome is the log equity premium, $\log(1+R_t)-\log(1+R_{f,t})$, with CRSP's S\&P~500 return. The predictors, dated $t-1$, are the log dividend--price ratio (dp), the log earnings--price ratio (ep), the term spread, long-term yield less T-bill rate (tms), and the default spread, BAA less AAA yield (dfy). All regressions and the historical mean start with 1927. The first forecast is for 1947, after 20 observations. The predictors, years, rows and transformations were fixed before any result was computed. Forecasts in percent:
\begin{center}\small
\begin{tabular}{@{}lrrrrrrr@{}}
\toprule
Year & Outcome & Hist.\ mean & dp & ep & tms & dfy & Mean comb. \\
\midrule
1947 & 4.31 & 4.10 & 3.04 & 3.87 & 3.73 & 3.62 & 3.57 \\
1970 & $-2.85$ & 6.98 & 3.42 & 4.14 & 1.00 & 7.29 & 3.96 \\
1990 & $-10.83$ & 5.86 & 0.51 & 3.28 & 3.81 & 5.86 & 3.37 \\
2010 & 13.99 & 5.47 & 0.29 & 2.11 & 11.81 & 5.41 & 4.90 \\
2025 & 12.47 & 6.46 & 2.08 & 3.08 & 4.38 & 6.14 & 3.92 \\
\bottomrule
\end{tabular}
\end{center}
Each recursive regression uses outcomes through $t-1$ and predictor rows dated $t-1$ for the forecast of year $t$; this is chronological indexing. The exercise does not model publication lags or data revisions (especially for E12), so these rows do not establish point-in-time availability or a real-time investable backtest.

\emph{Evaluation against the historical mean}, $R^2_{OS}$ in percent:
\begin{center}\small
\begin{tabular}{@{}lrrrrr@{}}
\toprule
Forecasts & dp & ep & tms & dfy & Mean comb. \\
\midrule
1947--2025 (79) & $-1.83$ & $-2.28$ & $-3.60$ & $-0.83$ & 2.56 \\
1947--2005 (59) & 1.34 & 4.61 & $-3.70$ & $-1.00$ & 5.82 \\
\bottomrule
\end{tabular}
\end{center}
Over 1947--2025 no single predictor beats the mean, but the combination does; its $\Delta$RMSE is 0.21 points. The Clark--West $t$ statistics of the combination, 1.44 and 2.13, are descriptive only, for the reason given in Conditions.

The four forecast errors are correlated between 0.89 and 0.98. Their deviations from the historical mean are correlated between $-0.30$ and 0.82, with dp and ep the one closely related pair. The deviations' standard deviations are 3.13, 5.84, 3.75 and 0.67 points, but that of their average is 2.06. In the decomposition above (all $\times100$), $\mathrm{MSE}_c=2.6500$ equals the average member's 2.7777 less the spread 0.1277; the historical mean's is 2.7197.

This is an annual illustration with four predictors, not a reproduction of the working paper's monthly Table~1. One window and one predictor set prove little, and a combination need not beat the mean in other samples.

\qffield{Pitfalls} Choosing the predictors, weights or evaluation start after seeing the results; with many tries, a data-mining correction is needed. Ex post optimal weights: from the 1947--2025 errors they are $-0.02$, 0.51, 0.40 and 0.12, which lowers the MSE only to 2.6380 and could not have been known in 1947. Reading a combination's gain as evidence that any one predictor works. Treating a Clark--West $p$-value for a combination as exact. Expecting $N$ to remove error variance when the errors share the return shock. Treating an average of point forecasts as a combined predictive distribution (\hyperref[meth:proper-scoring-rules]{proper scoring}).

\qffield{See also} \hyperref[meth:nested-forecast-tests]{nested forecast tests} (the single-predictor comparisons and $R^2_{OS}$); \hyperref[meth:restricted-forecast-value]{restricted forecasts} (sign restrictions and utility gains); \hyperref[meth:alpha-signal-combination]{combining alpha signals}; \hyperref[meth:forecast-accuracy-evaluation]{chronological forecast evaluation}.

\end{qfmethod}

\begin{qfmethod}[Scoring predictive distributions and intervals with proper scoring rules: the logarithmic score, the CRPS and the interval score, their sign conventions and units, and why a point loss cannot rank forecasts with the same mean]\label{meth:proper-scoring-rules}
\qffield{Statement} \emph{Point and probabilistic forecasts.} A point forecast is a single number, scored by a loss on its error. A probabilistic forecast is a predictive distribution $P$, or some of its quantiles, such as an interval. A \emph{scoring rule} $S(P,x)$ rewards the forecaster with $S(P,x)$ when $x$ materialises. Gneiting and Raftery take scores as positively oriented rewards to be maximised \citep[JASA 102(477), 2007, p.~359]{wide-gneiting-raftery}. With $S(P,Q)$ the expected score under $Q$, the rule is \emph{proper} if $S(Q,Q)\ge S(P,Q)$ for all $P$ and $Q$. It is \emph{strictly proper} if equality holds only for $P=Q$ (eq.~(1), p.~360). A forecaster who maximises expected score then has no reason to report anything but their true belief.

\emph{Squared error is not enough.} The negative squared error of the predictive mean, $-\|\mu_P-x\|^2$, is proper but not strictly proper (eq.~(23), p.~367). It depends on the mean alone, so forecasts with the same mean but different spreads tie.

\emph{Three rules} (each also given here as a negatively oriented loss, smaller is better):
\begin{itemize}
\item \emph{Logarithmic score}, $\mathrm{LogS}(p,x)=\log p(x)$ for a predictive density $p$ (eq.~(19), p.~365); loss $-\log p(x)$. Strictly proper among densities; its expected-score divergence is the Kullback--Leibler divergence.
\item \emph{Continuous ranked probability score}, $\mathrm{CRPS}(F,x)=-\int_{-\infty}^{\infty}\bigl(F(y)-\mathbb 1\{y\ge x\}\bigr)^2dy=\tfrac12E|X-X'|-E|X-x|$, with $X,X'$ independent draws from $F$ (eqs.~(20)--(21), p.~367). For $F=N(\mu,\sigma^2)$ and $z=(x-\mu)/\sigma$, the loss is $\sigma\bigl[z(2\Phi(z)-1)+2\varphi(z)-1/\sqrt\pi\bigr]$. Strictly proper among distributions with a finite mean.
\item \emph{Interval score} for a central $(1-\alpha)$ interval $[l,u]$, negatively oriented as printed (eq.~(43), p.~370):
\[ S^{\mathrm{int}}_\alpha(l,u;x)=(u-l)+\tfrac2\alpha(l-x)\mathbb 1\{x<l\}+\tfrac2\alpha(x-u)\mathbb 1\{x>u\}. \]
It charges the width and, beyond the interval, a penalty that grows as $\alpha$ falls. It equals $2/\alpha$ times the sum of the quantile (``tick'' or ``check'') losses $(\mathbb 1\{x\le r\}-\tau)(r-x)$ at $\tau=\alpha/2$ and $1-\alpha/2$ (eqs.~(41)--(42), p.~370).
\end{itemize}

\qffield{Conditions} \emph{Units.} The CRPS and the interval score are in the units of the outcome; for a point forecast the CRPS reduces to the absolute error, so point and probabilistic forecasts can be compared on one scale (p.~367). The log score is in nats and depends on the units of the density: a density of returns in percent is one hundredth of the density in decimals, so $-\log p$ is higher by $\log100=4.61$. Differences between forecasts of the same outcome are unit-free (Interpretation).

\emph{Behaviour.} The log score is local: it uses only the density at the outcome, and it punishes low-probability outcomes harshly. Gneiting and Raftery found it ``highly sensitive to extreme cases'' and the CRPS less so (p.~374). The linear score $p(x)$ looks natural but is improper: it rewards a density concentrated at the mode (p.~366).

\emph{Propriety is about expectations.} A strictly proper score rewards the true distribution on average. Over a finite sample, a wrong forecast can score better, so a difference in average scores needs its own test of equal predictive ability (Interpretation).

\emph{Combining densities.} Averaging point forecasts and combining predictive distributions are different operations. The equal-weight \emph{linear pool} of two densities is their mixture. Its variance adds the squared spread of the means to the average variance. Choosing a normal density centred on the averaged point forecast and assigning it the average component variance is an additional modelling assumption: point forecasts alone determine neither a variance nor a predictive distribution (Interpretation).

\qffield{Practice} \emph{Same mean, different spread.} Two forecasts of an annual excess return, in percentage points: $A=N(5,10^2)$ and $B=N(5,20^2)$, with 90 percent central intervals $[-11.45,21.45]$ and $[-27.90,37.90]$. Losses (smaller is better):
\begin{center}\small
\begin{tabular}{@{}rrrrrrrr@{}}
\toprule
& Squared error & \multicolumn{2}{c}{$-\log p(x)$} & \multicolumn{2}{c}{CRPS} & \multicolumn{2}{c}{Interval score} \\
$x$ & $A$ and $B$ & $A$ & $B$ & $A$ & $B$ & $A$ & $B$ \\
\midrule
5 & 0 & 3.222 & 3.915 & 2.337 & 4.674 & 32.90 & 65.79 \\
$-15$ & 400 & 5.222 & 4.415 & 14.528 & 12.049 & 103.93 & 65.79 \\
$-40$ & 2025 & 13.347 & 6.446 & 39.358 & 33.886 & 603.93 & 307.85 \\
\bottomrule
\end{tabular}
\end{center}
The squared error cannot separate $A$ and $B$. The three proper scores favour the sharper $A$ when the outcome is at the mean and the wider $B$ after a loss of 20 or 45 points. With $x$ in decimals, every $-\log p$ falls by 4.605; at $x=5$ the log loss of $A$ is $-1.384$, a negative loss. The CRPS values agree with numerical integration of eq.~(20) and with simulation of eq.~(21).

\emph{Propriety, numerically.} If returns are truly $N(5,15^2)$, forecasts $N(5,s^2)$ with $s=10$, 15 and 20 have expected log losses of 4.347, 4.127 and 4.196. Their expected CRPS is 8.742, 8.463 and 8.663, and their expected 90 percent interval scores 74.35, 61.88 and 68.78. Each loss is smallest at the true $s=15$ on a grid from 5 to 30. The expected linear score instead keeps rising as $s$ shrinks.

\emph{Interval forecasts for a heteroskedastic process} \citep[eqs.~(44)--(47), Table~2, p.~371]{wide-gneiting-raftery}. For $X_{t+1}=\tfrac12X_t+\tfrac12X_t\varepsilon_t+\varepsilon_t$ with standard normal $\varepsilon_t$, three one-step 95 percent intervals are compared: $I$ from the true conditional distribution, $J$ from the unconditional quantiles, and $K$, which minimises expected width at nominal coverage. Over 100,000 forecasts:
\begin{center}\small
\begin{tabular}{@{}lrrr@{}}
\toprule
Interval & Coverage (\%) & Average width & Average interval score \\
\midrule
$I$ & 94.98 (95.01) & 4.00 (4.00) & 4.76 (4.77) \\
$J$ & 94.91 (95.08) & 5.46 (5.45) & 8.03 (8.04) \\
$K$ & 94.94 (94.98) & 3.79 (3.79) & 5.29 (5.32) \\
\bottomrule
\end{tabular}
\end{center}
Published values are in brackets. All three intervals have nominal coverage and $K$ is the narrowest, but the conditional interval $I$ scores best: $K$ collapses to a point when the conditional variance is highest. Coverage alone cannot rank the three.

\emph{A linear pool versus an averaged point forecast.} Forecasts $N(0,10^2)$ and $N(8,10^2)$ average to the point forecast 4. Their equal-weight pool has mean 4 and variance $100+4^2=116$; the normal density around the averaged forecast, $N(4,10^2)$, has variance 100. At $x=4$, $N(4,10^2)$ has the lower log loss and CRPS (3.222 and 2.337 against the pool's 3.302 and 2.527). At $x=-15$ and $x=20$ the pool does (log losses 4.842 and 4.389 against 5.027 and 4.502). If the pool is the true distribution, its expected log loss is 3.796 against 3.802, and its expected CRPS 6.082 against 6.090, as strict propriety requires.

All forecasts in this field are Gaussian, or mixtures of Gaussians, by assumption. Return distributions have fatter tails, which raises the log loss of thin-tailed forecasts in crashes.

\qffield{Pitfalls} Ranking density forecasts by the squared error of their means. Comparing log scores computed in different units, or reading a negative log loss as an error. Mixing positively and negatively oriented scores in one table. Judging intervals by coverage alone. Treating an average of point forecasts, with the members' own spread, as the combined predictive distribution. Using an improper score such as the linear score, which favours overconfident forecasts.

\qffield{See also} \hyperref[meth:forecast-combination]{forecast combination} (averaging point forecasts); \hyperref[meth:forecast-accuracy-evaluation]{chronological forecast evaluation} (RMSE favours the mean, MAE the median); \hyperref[meth:volatility-forecast-loss]{volatility forecast losses}; \hyperref[meth:var-es-estimation-methods]{VaR and ES estimation} (VaR is a quantile forecast); \hyperref[meth:conditional-predictive-ability]{conditional predictive ability} (testing differences in scores).

\end{qfmethod}

\begin{qfmethod}[Testing equal conditional predictive ability of forecasting methods: the null on losses at the estimates used, instruments known at the forecast origin, the Wald statistic and its covariance, bounded estimation windows, and when the conditional null can hold]\label{meth:conditional-predictive-ability}
\qffield{Statement} \emph{What is compared.} A \emph{forecasting method} is a model together with its estimation procedure and the data used to estimate it, including the length of the window. The null concerns losses at the estimates actually used, conditional on an information set $\mathcal G_t$:
\[ H_0:\ E\bigl[L(y_{t+\tau},\hat f_t)-L(y_{t+\tau},\hat g_t)\,\big|\,\mathcal G_t\bigr]\equiv E[\Delta L_{t+\tau}\mid\mathcal G_t]=0 \]
\citep[authors' version, April 2006, eqs.~(2)--(3), pp.~5, 8]{giacomini-white-2006-wp}. The Diebold--Mariano--West null (eq.~(1)) evaluates the losses at the parameters' probability limits and so compares \emph{models}. The method null counts estimation noise, so a small misspecified model can be as accurate as a large correct one (pp.~5--6).

\emph{Conditional and unconditional.} With $\mathcal G_t$ trivial, the null is equal average loss: which method was better on average, the question for an unspecified future date. With $\mathcal G_t=\mathcal F_t$, the information at the forecast origin, it says that nothing known then predicts which method will be more accurate at $t+\tau$ (pp.~3, 6). The conditional null implies the unconditional one, not the reverse.

\emph{Instruments.} For $\tau=1$ and $\mathcal G_t=\mathcal F_t$ the null makes $\Delta L_{t+1}$ a martingale difference sequence, so $E[h_t\Delta L_{t+1}]=0$ for every $q\times1$ vector $h_t$ known at $t$ (the ``test function''). The test checks only this implication for the chosen $h_t$, such as a constant and the lagged loss difference, or a business-cycle indicator. It can ``accept'' a false null when $\Delta L$ is correlated with information left out of $h_t$, and weak instruments dilute its power (pp.~9, 11).

\emph{Statistic.} With $Z_{t+1}=h_t\Delta L_{t+1}$ over the $n$ forecasts,
\[ T=n\,\bar Z'\hat\Omega^{-1}\bar Z,\qquad \hat\Omega=n^{-1}\textstyle\sum_t Z_{t+1}Z_{t+1}', \]
and the null is rejected when $T$ exceeds the $1-\alpha$ quantile of $\chi^2_q$ (eq.~(4), p.~9). $\hat\Omega$ is uncentred and omits autocovariances: under the null $Z$ is a martingale difference. $T$ equals $n$ times the uncentred $R^2$ of a regression of a constant one on $Z'$. A HAC estimator is also valid but may lose power (comments~5 and~7, p.~10).

\emph{When serial correlation must be allowed for.} For $\tau>1$ the conditional null leaves $h_t\Delta L_{t+\tau}$ correlated up to lag $\tau-1$, and $\tilde\Omega$ adds those autocovariances with weights tending to one (eq.~(6), Theorem~3, pp.~11--12). The unconditional null implies no dependence structure: its $t$ statistic, $\overline{\Delta L}/(\hat\sigma/\sqrt n)$, needs a HAC variance with a truncation lag growing with $n$, and it coincides with Diebold and Mariano's (eq.~(8), Theorem~4, pp.~12--13). Rolling windows make $\Delta L$ itself dependent. Zhu and Timmermann therefore replace the martingale-difference variance with Newey--West, or with a subsample $t$: $K$ block means against Student $t_{K-1}$ \citep[arXiv:2006.03238v2, Remarks~1--2 and eq.~(13), p.~14]{zhu-timmermann-2022-arxiv}.

\qffield{Conditions} \emph{Assumptions of the one-step test} \citep[Theorem~1, p.~9]{giacomini-white-2006-wp}: the estimation window is bounded, $m\le\bar m<\infty$, while $n\to\infty$, so estimation noise does not vanish; data and instruments are mixing, allowing heterogeneity and breaks at unknown dates; moments of order $2(r+\delta)$ exist; and $\Omega_n$ is uniformly positive definite. The degrees of freedom are the number of linearly independent instruments. Nested models are allowed, because non-vanishing estimation noise prevents the singularity that the model null meets (comment~4, p.~10). Any loss qualifies (squared, absolute, linex, direction of change, predictive log likelihood; p.~8), so a \hyperref[meth:proper-scoring-rules]{proper score} can define $\Delta L$.

\emph{Singular instrument covariance.} $\hat\Omega=n^{-1}Z'Z$, with rows $Z_{t+1}'=h_t'\Delta L_{t+1}$, is singular exactly when the $n\times q$ matrix $Z$ has rank below $q$. A single constant, $h_t=1$, is valid: then $\hat\Omega$ is the mean of $\Delta L^2$, positive unless every $\Delta L$ is zero. Singularity arises when the instruments are linearly dependent over the sample, for example a second constant, or an indicator that never changes, beside the intercept; $T$ is then undefined. A generalised inverse returns the statistic of the reduced set, whose degrees of freedom are the rank, not $q$ (Interpretation; checked numerically). Drop the redundant instrument and report the $q$ used.

\emph{Expanding windows are outside the theory.} A bounded window ``rules out an expanding window forecasting scheme'' \citep[p.~8 and comment~3, p.~10]{giacomini-white-2006-wp}; a single fixed estimation sample is allowed. The recursive comparisons under \hyperref[meth:nested-forecast-tests]{nested forecast tests} and \hyperref[meth:forecast-combination]{forecast combination} are therefore not Giacomini--White exercises. Their methods include an expanding window, and Clark--West's critical values address the population-slope null with $R,P\to\infty$. Recasting them needs a fixed window and new forecasts, as in the illustration below. With an expanding window the statistic's limit can be non-standard \citep[Proposition~8, p.~17]{zhu-timmermann-2022-arxiv}.

\emph{When the conditional null can hold.} Under squared loss $\Delta L=(f_1+f_2-2y)(f_1-f_2)$. So, if $f_1\neq f_2$, $E[\Delta L_{t+1}\mid\mathcal F_t]=0$ exactly when $E[y_{t+1}\mid\mathcal F_t]=(f_{1,t}+f_{2,t})/2$ \citep[Proposition~1, p.~5]{zhu-timmermann-2022-arxiv}. With rolling-window forecasts it fails whenever the conditional mean depends on data older than the window (Proposition~2, p.~6). For an autoregression with estimated coefficients and a window longer than its order it should not be expected (Lemma~1, p.~8). For nested models whose average losses are equal by construction, $\Delta L$ is still not a martingale difference (Proposition~3, p.~8). So equal population models do not establish equal method accuracy, and equal average accuracy does not establish the conditional null. A rejection of the conditional null therefore does not by itself show that either method is better on average.

\emph{Versions.} Giacomini and White are cited from the authors' post-referee version of April 2006. The \emph{Econometrica} article (74(6), 2006, 1545--1578; catalogue \nolinkurl{wide-giacomini-white}) was not obtained, and its numbering may differ: the copy's text calls its tests ``Theorems~5 and~6''. Zhu and Timmermann are cited from arXiv version~2 (August 2022). Version~1 (June 2020) claimed instead that the null cannot be valid under a rolling window.

\qffield{Practice} \emph{Same average accuracy, different conditional accuracy} (constructed). A state $S_t$, boom or bust with probability $\tfrac12$, is known at $t$. Then $y_{t+1}=1+\varepsilon$ after a boom and $-1+\varepsilon$ after a bust, with $\varepsilon\sim N(0,1)$. Method~A always forecasts $+1$, method~B $-1$. Expected squared losses are 1 and 5 for A, and 5 and 1 for B, so both average 3. $\Delta L=L_A-L_B=-4y_{t+1}$ has conditional mean $-4$ after a boom and $+4$ after a bust. Neither method is better unconditionally; conditionally, A is better after a boom. With $n=120$ and 10,000 samples, the unconditional test rejects 5.2 percent of the time (Monte Carlo s.e.\ 0.2 points), the conditional test with $h_t=(1,S_t)$ every time.

\emph{Simulated methods, not simulated loss differences.} Each design simulates the data and both forecasting methods. The status of each null is derived from them. All tests are at 5 percent, with 10,000 samples except C (2,000); Monte Carlo standard errors are at most 0.004.
\begin{center}\small
\begin{tabular}{@{}llrr@{}}
\toprule
Design and derived null & Test & $n=120$ & $n=1000$ \\
\midrule
N: $y=\tfrac12(f_1+f_2)+\varepsilon$, conditional null true & $h=(1,\Delta L_t)$ & 0.042 & 0.047 \\
U: unconditional null true, & $h=1$, martingale-difference variance & 0.026 & 0.024 \\
\quad conditional null false & $t$, uniform HAC, lag 4 & 0.093 & 0.055 \\
 & subsample $t$, $K=4$ & 0.050 & 0.046 \\
 & $h=(1,f_1)$ & 0.061 & 0.044 \\
 & $h=(1,f_1^2)$ & 0.142 & 0.788 \\
C: equal population models, method null false & $h=(1,\Delta L_t)$ & & 0.999 \\
\bottomrule
\end{tabular}
\end{center}
In N each forecast is the 12-period moving average of its own AR(1) signal. The null holds only because the outcome is defined as their average (Zhu and Timmermann's Example~2). In U, $y=\tfrac12+\varepsilon$, $f_1$ is the mean of the last $m=4$ outcomes and $f_2=0$ (their Example~4). Then $E[\Delta L_{t+1}\mid\mathcal F_t]=v_t^2-\tfrac14$, with $v_t=f_{1,t}-\tfrac12$, so $E[f_1\Delta L]=0$ but $E[f_1^2\Delta L]=1/8$: the instrument $f_1$ cannot see the failure and $f_1^2$ can (Interpretation). $\Delta L$ is $m$-dependent, so the unconditional null needs autocovariances. Without them, the variance of the normalised statistic is $(2m+1)/(3m)=0.75$ for Gaussian errors (derived here, Interpretation; simulated 0.755). Zhu and Timmermann's Proposition~6 gives 0.677 (p.~14); its proof's Gaussian-case fourth moment omits the factor 3 on the pairings (eq.~(23), p.~25, and before eq.~(28), p.~26), and with it the same steps give 0.75 (Interpretation). Both make the test undersized. In C both methods use 20 observations: the window mean, and OLS on an AR(1) predictor whose true slope is zero. Then $E[L_M-L_D]=-E[\hat b^2(x_t-\bar x)^2]<0$ (Interpretation; simulated $-0.097$, s.e.\ 0.0004). The test rejects, and the sign of $\overline{\Delta L}$ favours the window mean.

For Zhu and Timmermann's expanding-window Example~5, with $T=2000$, comparing the statistic with 1.96 rejects 5.7 and 5.4 percent of the time ($\lambda=0.25$, 0.5). That agrees with the limit derived here, 5.4 and 5.3 percent. It does not reproduce their Table~2, 19.6 and 15.8 percent (p.~18). Their eq.~(14) has the opposite sign on the stochastic-integral term to this derivation and to their own proof (eq.~(30), p.~28); simulated as printed, it gives 14.9 and 11.7 percent. Neither sign gives Table~2, whose values remain unexplained here.

\emph{Annual equity premium with a fixed 20-year window.} The window, methods, losses, instrument and samples were fixed before any result. The outcome is the annual log equity premium in percent. Method~M is the mean of the previous 20 outcomes; D is OLS on the lagged log dividend--price ratio over the same 20 years. The 1947 forecasts (4.10 and 3.04 percent) equal the recursive ones of the earlier entries; later the windows differ. $\Delta L=L_M-L_D$, so a positive mean favours D. The CRPS assumes a Gaussian $N(\hat f,s^2)$, $s^2$ being the window's variance (M) or residual variance (D).
\begin{center}\small
\begin{tabular}{@{}llrrrrr@{}}
\toprule
Forecasts & Loss & Mean $\Delta L$ & DM $t$ & $p$ & $T$, $h=(1,\Delta L_t)$ & $p$ \\
\midrule
1947--2025 (79) & squared (pct$^2$) & 9.87 & 0.50 & 0.62 & 2.73 & 0.26 \\
 & CRPS (pct) & 0.116 & 0.31 & 0.75 & 3.01 & 0.22 \\
1947--2005 (59) & squared (pct$^2$) & 3.38 & 0.14 & 0.89 & 2.63 & 0.27 \\
 & CRPS (pct) & $-0.001$ & $-0.00$ & 1.00 & 2.76 & 0.25 \\
\bottomrule
\end{tabular}
\end{center}
The DM $t$ uses Newey--West with lags 3 and 2; $T$ uses 78 and 58 observations and $\chi^2_2$. D's $R^2_{OS}$ against M is 3.40 and 1.22 percent. In the recursive comparison of the earlier entries, dp's $R^2_{OS}$ against the historical mean for 1947--2025 is $-1.83$ percent. Fixing the window changes both methods, so the two numbers answer different questions. No test rejects, conditionally or unconditionally, and with 78 observations the $\chi^2$ approximation is rough. The predictor years are calendar-indexed, from the 2026 vintage of Goyal's file, with no publication lags or revisions: this is not a point-in-time backtest.

\qffield{Pitfalls} Reporting a Giacomini--White test on expanding-window forecasts. Reading a rejection of the conditional null as evidence that one method is better on average, or a non-rejection with a small $h_t$ as evidence of equal conditional accuracy. Using the martingale-difference variance for the unconditional null or for $\tau>1$. Choosing the window, instruments or evaluation sample after seeing results. Keeping $q$ degrees of freedom after replacing a singular $\hat\Omega$ with a generalised inverse. Reading a Clark--West result, which concerns the population slope, as one about forecasting methods.

\qffield{See also} \hyperref[meth:nested-forecast-tests]{nested forecast tests} (the model null and recursive windows); \hyperref[meth:forecast-combination]{forecast combination}; \hyperref[meth:proper-scoring-rules]{proper scoring rules} (losses for density and interval forecasts); \hyperref[meth:forecast-accuracy-evaluation]{chronological forecast evaluation}; \hyperref[meth:correlation-independence-tests]{overlapping observations}; \hyperref[meth:hypothesis-testing]{hypothesis tests}.

\end{qfmethod}
